\documentclass[11pt,a4paper]{article}

% Compile with pdfLaTeX (the Overleaf default). Upload the figures/ folder too.
\usepackage[utf8]{inputenc}
\usepackage[T1]{fontenc}
\usepackage[english]{babel}
\usepackage{lmodern}
\usepackage{microtype}
\usepackage[a4paper,margin=2.3cm]{geometry}
\usepackage{amsmath}
\usepackage{graphicx}
\usepackage{booktabs}
\usepackage{tabularx}
\usepackage{array}
\usepackage{float}
\usepackage{caption}
\usepackage{xcolor}
\usepackage{listings}
\usepackage[hidelinks]{hyperref}

\graphicspath{{figures/}}
\captionsetup{font=small,labelfont=bf}
\setlength{\parskip}{5pt}
\setlength{\parindent}{0pt}
% Long names in typewriter type (RandomNumberGenerator, "proto":"lanchat/1") cannot be hyphenated:
% let TeX stretch the spaces of such a line a little instead of letting the name stick out.
\setlength{\emergencystretch}{3em}
\renewcommand{\arraystretch}{1.15}

\definecolor{wrong}{HTML}{B42828}
\definecolor{notebg}{HTML}{F3F2EE}
% Red marker for every value that still has to be filled in from your own results.
% Search the file for TODO before handing the report in.
\newcommand{\TODO}[1]{\textcolor{wrong}{\textbf{[TODO: #1]}}}
% Include a figure if the file exists, otherwise show a visible placeholder.
\newcommand{\fig}[2]{\IfFileExists{figures/#1}{\includegraphics[width=#2]{#1}}%
  {\fbox{\parbox{0.9\linewidth}{\centering\vspace{2em}\TODO{figure \detokenize{#1} not generated yet}\vspace{2em}}}}}
\newcommand{\question}[1]{\par\medskip\textbf{#1}\par}
\newcommand{\notebox}[1]{\par\smallskip\noindent\colorbox{notebg}{\parbox{\dimexpr\linewidth-2\fboxsep}{\small #1}}\par\smallskip}
% Code and program output: no extra space after a full stop or colon inside it.
\newcommand{\code}[1]{{\frenchspacing\texttt{#1}}}

% Footnotesize (not small as in lab 1): a log line or a datagram fits on one line.
\lstset{basicstyle=\ttfamily\footnotesize,columns=fullflexible,keepspaces=true,
        frame=single,rulecolor=\color{black!25},backgroundcolor=\color{notebg},
        xleftmargin=3pt,xrightmargin=3pt,breaklines=true}

\begin{document}

% Title
\begin{center}
  {\small NETWORK PROGRAMMING \textperiodcentered{} UTM}\\[6pt]
  {\LARGE\bfseries Laboratory Work 3: A Local Network Chat}\\[10pt]
  \begin{tabular}{rl}
    Authors: & Pancenco Ina, Rusnac Nichita \\
    Group:   & FAF-242 \\
    Teacher: & Ernest B\^{\i}tca \\
    Date:    & November 2026
  \end{tabular}
\end{center}
\vspace{4pt}\hrule\vspace{8pt}

% Section 1
\section{The task and the environment}

The goal of this laboratory work is a chat for the local network built directly on UDP sockets.
There is no server: every running copy of the program is a peer that both sends and receives. A
message of the general channel reaches every peer with a single datagram, sent to a multicast group
(or as a broadcast); a private message goes to one peer only and, in the last task, is encrypted
with a key derived from a passphrase. The peers find each other without a server, notice who joins
and who leaves, and never crash because of a datagram that is not a message of the chat.

The chat is written in C\# as part of the same \code{Labs} program as the previous laboratory works
(\code{Labs.exe lab3 chat -{}-name Ana -{}-port 5000}). For the bonus, a second implementation of the
same protocol was written in Java, and the two talk to each other. The format of the datagrams is
specified in \code{lab3/PROTOCOL.md}, so that a third program could join the same chat.

Only the socket API is used: \code{Socket} with \code{SocketType.Dgram} in C\# and
\code{DatagramSocket} in Java, with the equivalents of \code{socket}, \code{bind},
\code{setsockopt}, \code{sendto}, \code{recvfrom}, \code{select} and \code{close}. No
\code{UdpClient}, no library for messaging or discovery. The JSON of the datagrams is written and
read with \code{System.Text.Json} in C\#; Java has no JSON in its standard library, so the Java peer
has its own small reader and writer (\code{Json.java}, 210 lines).

\begin{center}
\small
\begin{tabularx}{\linewidth}{@{}lX@{}}
\toprule
Component & Value \\
\midrule
Computer & laptop, Intel Core i7-13700H (14 cores, 20 logical processors), 32\,GB RAM \\
Operating system & Windows 10 Home 22H2 (build 19045) \\
Network & Wi-Fi (MediaTek Wi-Fi 6E MT7902), address 192.168.100.15/24 in a home network;
  the computer also has two link-local virtual adapters (169.254.x.x) and the loopback adapter \\
C\# & .NET 9 (SDK 9.0.304, runtime 9.0.8), Release build \\
Java & JDK 26.0.1, compiled with \code{-{}-release 17}, so any JDK from 17 on runs it \\
Tests and charts & Python 3.13, standard library only (charts: matplotlib) \\
\bottomrule
\end{tabularx}
\end{center}

All the tests and measurements of this report ran on this one computer, with up to sixteen peers
at the same time. The general channel still left the computer through the Wi-Fi adapter: the system
sends a multicast or broadcast datagram on the network and gives a copy of it to every local socket
of the group. What a single computer cannot show is a datagram that arrives from another computer;
that is shown at the defense, as the assignment requires (Section~\ref{sec:two}).

% Section 2
\section{Design}

\subsection{Two sockets per peer}

The central decision follows from a property of shared ports. Several peers on one computer must
receive the general channel, so they all bind the port of the chat; but a unicast datagram sent to a
port that several sockets share reaches only one of them (Section~\ref{sec:task1}, question~2).
One socket cannot do both jobs, so every peer has two:

\begin{center}
\small
\begin{tabularx}{\linewidth}{@{}lp{2.3cm}>{\raggedright\arraybackslash}p{4.9cm}X@{}}
\toprule
Socket & Bound to & Options & What it receives \\
\midrule
shared & \code{0.0.0.0:5000} & \code{SO\_REUSEADDR} before \code{bind}; joins the group &
  the general channel: \code{hello}, \code{join}, \code{alive}, \code{bye}, \code{chat} \\
own & \code{0.0.0.0:0} & \code{IP\_MULTICAST\_IF}, \code{IP\_MULTICAST\_TTL}\,=\,1, \code{SO\_BROADCAST} &
  only what is for this peer: \code{here}, \code{taken}, \code{private}, \code{ack} \\
\bottomrule
\end{tabularx}
\end{center}

Port 0 asks the system for any free port, so the own socket has a port that no other peer uses,
for example 56437. \textbf{Every datagram of a peer is sent from its own socket}, also the ones to
the group. Its source address and port are therefore the address of the peer in the chat: the
others learn it from any datagram, and send their answers and private messages there. No peer ever
has to tell its address in the content of a message.

\subsection{Threads}

The C\# peer has two threads. The main thread reads the keyboard and sends the general and private
messages. The network thread waits for both sockets at once with \code{Socket.Select}, which returns
as soon as a datagram waits on either of them, or after 0.2\,s. It receives and handles every
datagram, and on every pass also checks the timers: send \code{alive} every 5\,s, remove the peers
not heard from for 15\,s, and send again the private messages without an acknowledgement.

The data the two threads share is protected as in laboratory work~2: the table of the peers online,
the private messages waiting for an acknowledgement, the state of the peer (joining, online,
refused, left), the sending on the own socket, the log file and the console each have a lock. The
table holds immutable records, which are replaced as a whole, so a reader never sees half a change.

Messages from the network must appear while the user is typing. A class \code{ConsoleUi} reads the
keyboard key by key and keeps the line being typed; when a message arrives, it erases that line,
prints the message, and draws the line again below it. The user's typing is never interrupted and
never mixed with a message. When the input is not a terminal (a pipe, as in the tests), it falls
back to ordinary line reading.

\subsection{The datagrams}

A datagram is one JSON object in UTF-8. UDP keeps the boundaries of datagrams, so one datagram is
always one complete message, with no delimiter and no length prefix. Every datagram carries the
version of the protocol, its kind, the name of the sender and the \emph{id} of the running sender:
8 random bytes, written as 16 hex digits, new at every start. A general message:

\begin{lstlisting}
{"proto":"lanchat/1","type":"chat","from":"Ana","id":"5f0c9a1e7b2d4c33","seq":4,"text":"hello everyone"}
\end{lstlisting}

\begin{center}
\small
\begin{tabularx}{\linewidth}{@{}llX@{}}
\toprule
Kind & Sent to & Meaning \\
\midrule
\code{hello} & general channel & a peer starts: ``I want to join with this name'' \\
\code{here}  & the new peer & answer to \code{hello}: ``I am online'' \\
\code{taken} & the new peer & answer to \code{hello}: ``this name is mine'' \\
\code{join}  & general channel & one second after \code{hello} without \code{taken}: ``I have joined'' \\
\code{alive} & general channel & every 5 seconds: ``I am still here'' \\
\code{bye}   & general channel & \code{/quit}, end of the input, Ctrl+C \\
\code{chat}  & general channel & a message for everybody (\code{seq}, \code{text}) \\
\code{private} & the receiver's own socket & a private message (\code{seq}, \code{to}, and \code{text} or \code{nonce} and \code{data}) \\
\code{ack}   & the sender's own socket & ``your private message with this \code{seq} arrived'' \\
\bottomrule
\end{tabularx}
\end{center}

The receiver checks every field: the datagram must be valid UTF-8 and JSON, have
\code{"proto":"lanchat/1"} and a known kind, a valid name (1 to 16 letters, digits, \code{\_} or
\code{-}), a valid id, and the fields of its kind with valid values (a text of 1 to 500 characters,
a sequence number of at least 1, a nonce of 12 bytes). A datagram longer than 4\,096 bytes is
dropped before it is even parsed. The receive buffer has 65\,536 bytes, so that a datagram that is
too long arrives whole and can be measured and dropped, instead of being silently cut.

\subsection{The log}

Every peer writes a log file with a line for every datagram that it sends, receives or drops: the
time, what happened, the endpoint of the other side, the kind, the size and a note. Lines starting
with \code{-{}-{}-{}-} record other events. The log of Ana in a session with Ion and Maria on port 5000:
she joins, Ion and Maria join, the three send a general message each, Ana sends a private message
to Ion, a test program sends a datagram that is not a message, and Ion leaves:

\begin{lstlisting}
00:23:18.191  sent  239.255.42.42:5000     hello       73 B
00:23:18.204  recv  192.168.100.15:60685   hello       73 B  own datagram, not shown again
00:23:19.196  sent  239.255.42.42:5000     join        72 B
00:23:19.196  recv  192.168.100.15:60685   join        72 B  own datagram, not shown again
00:23:19.340  recv  192.168.100.15:60686   hello       73 B
00:23:19.343  sent  192.168.100.15:60686   here        72 B
00:23:20.346  recv  192.168.100.15:60686   join        72 B
00:23:20.503  recv  192.168.100.15:63791   hello       75 B
00:23:20.505  sent  192.168.100.15:63791   here        72 B
00:23:21.516  recv  192.168.100.15:63791   join        74 B
00:23:21.528  sent  239.255.42.42:5000     chat       104 B
00:23:21.529  recv  192.168.100.15:60685   chat       104 B  own datagram, not shown again
00:23:21.538  recv  192.168.100.15:63791   chat        98 B
00:23:21.542  recv  192.168.100.15:60686   chat       104 B
00:23:21.545  sent  192.168.100.15:60686   private    119 B
00:23:21.547  recv  192.168.100.15:60686   ack         79 B
00:23:21.556  drop  192.168.100.15:63792   ?           24 B  not JSON
00:23:22.170  recv  192.168.100.15:60686   bye         71 B
\end{lstlisting}

Ana's own socket is \code{192.168.100.15:60685}, Ion's \code{:60686} and Maria's \code{:63791}: all
three run on this computer, and the source of every datagram tells who sent it.

% Section 3
\section{Task 1: The general channel}\label{sec:task1}

The shared socket of every peer is bound to \code{0.0.0.0} and the port of the chat, after
\code{SO\_REUSEADDR}, and joins the multicast group \code{239.255.42.42} on the network the chat
uses. A line typed by the user is sent with one \code{sendto} from the own socket to
\code{239.255.42.42:5000}, which is one datagram. The system sends it once on the Wi-Fi network and
gives a copy to every socket of this computer that is bound to port 5000 and has joined the group;
the network delivers it to the other computers whose sockets have joined it.

The group address lies in \code{239.255.0.0/16}, the \emph{local scope} of IPv4 multicast
(RFC~2365), meant for the network of one organisation. The time to live of our multicast datagrams
is 1, so no router passes them to another network. By default the chat uses the adapter that has a
default gateway, the real local network, and not one of the virtual adapters of the computer; the
option \code{-{}-interface} chooses another one, and \code{lab3 networks} lists them.

The output of Ion in the session of the log above (the first line, with the name of the Wi-Fi
adapter, is left out):

\begin{lstlisting}
This peer: 192.168.100.15:60686. Private messages: not encrypted (no --passphrase).
Log: D:\Second_year\PR-Labs\lab3\logs\tests\Ion-5000-2.log
Joining as Ion...
* You joined the chat as Ion. Online: Ana.
* Maria joined the chat
[general] Ana: hello everyone
[general] Maria: hi Ana
[general] Ion: hello from Ion
[private] Ana: see you at five
\end{lstlisting}

The test \code{general channel} checks that a message typed in one of three peers is shown by all
three, the author included, exactly once (it waits one more second for a copy that should not
come), and that the author's log has exactly one \code{chat} datagram sent, to the group.

\subsection{Answers to the questions of Task 1}

\question{1. Why is the socket bound to all the addresses of the computer (0.0.0.0), and not to its
IP address or to 127.0.0.1?}

The address in \code{bind} is a filter on the \emph{destination} address of the datagrams the socket
accepts. The datagrams of the general channel are not addressed to the IP address of the computer
at all: their destination is \code{239.255.42.42}, or \code{192.168.100.255} with
\code{-{}-broadcast}. \code{0.0.0.0} accepts every destination that arrives at this computer: the
group, the broadcast address, and any address of its adapters. A socket bound to
\code{127.0.0.1} only accepts datagrams sent to the loopback address, which never come from
another computer: it would hear nothing of the chat. A socket bound to the computer's own address
depends on the system. A small experiment on this computer bound one socket to each address, all to
the same port, and sent three datagrams of each kind:

\begin{center}
\footnotesize
\begin{tabular}{lcccc}
\toprule
Sent to \textbackslash{} bound to & \code{0.0.0.0} (1st) & \code{0.0.0.0} (2nd) & \code{192.168.100.15} & \code{127.0.0.1} \\
\midrule
the group \code{239.255.42.42}     & 3/3 & 3/3 & 3/3 & 0/3 \\
broadcast \code{192.168.100.255}   & 3/3 & 3/3 & 3/3 & 0/3 \\
broadcast \code{255.255.255.255}   & 3/3 & 3/3 & 3/3 & 0/3 \\
unicast \code{192.168.100.15}      & 0/3 & 0/3 & 3/3 & 0/3 \\
unicast \code{127.0.0.1}           & 0/3 & 0/3 & 0/3 & 3/3 \\
\bottomrule
\end{tabular}
\end{center}

Windows delivers multicast and broadcast datagrams also to a socket bound to the address of the
adapter on which they arrive; Linux does not, because their destination is not that address.
Even on Windows, a socket bound to the computer's address ties the program to one adapter and to
an address that DHCP may change while the chat runs. \code{0.0.0.0} works on every system and with
any address. The adapter is chosen elsewhere: when the socket joins the group
(\code{IP\_ADD\_MEMBERSHIP} with the address of the adapter) and for sending
(\code{IP\_MULTICAST\_IF}). The second one matters on this computer: a datagram sent to the group
from a PowerShell socket without \code{IP\_MULTICAST\_IF} left through one of the virtual adapters
and never reached the chat. The same bytes sent to a peer's own endpoint arrived and were dropped
and logged, and the test program, which sets the option, reaches the group every time.

\question{2. How can several peers on the same computer bind the same port? Which of them receive a
message of the general channel?}

Normally a second \code{bind} to a port in use fails with ``address already in use''. Every peer
sets \code{SO\_REUSEADDR} on its shared socket before \code{bind} (in C\#
\code{SocketOptionName.ReuseAddress}, together with \code{ExclusiveAddressUse = false}); the system
then lets every socket that asked for it bind the same port. A multicast or broadcast datagram is
delivered to \emph{all} of them: the system makes one copy for every socket bound to the port (and,
for multicast, member of the group). The table above shows it: both \code{0.0.0.0} sockets got
every multicast and broadcast datagram. So all the peers of one computer receive the general
channel, each with its own copy.

A unicast datagram, by contrast, is delivered to \emph{one} of the sockets only. With three sockets
bound to \code{0.0.0.0:47556}, ten unicast datagrams all went to the socket bound first; the other
two received nothing. Windows does not promise which one gets them. This is why nothing for one
peer only is ever sent to the shared port, and why every peer has its own socket.

\question{3. Your own message comes back to you. How does your chat recognise it?}

The datagram goes to the group, and the shared socket of the author is a member of the group like
any other, so the author receives its own message too. Every datagram carries the \emph{id} of its
sender: 8 random bytes chosen when the program starts. The network thread compares the id of every
datagram with its own; when they are equal, the datagram is its own and is not shown (the log says
\code{own datagram, not shown again}). The author's message is shown once, at the moment it is
sent, so the author sees it even if the echo is lost.

The id is used rather than the name, because a second program may try to use the same name (Task~2),
and rather than the source address, because it is the same in every language and does not depend on
which address the system chose. Besides, every message is shown only once: the peer remembers the
pair (id, sequence number) of the last 1\,000 messages and drops a copy, for example a message that
arrives twice through two networks.

\question{4. Why at most 500 characters? What would happen to a message of 100\,000 characters?}

A message must fit in one datagram, and a datagram should fit in one Ethernet frame, whose payload
for UDP is 1\,472 bytes (1\,500 bytes of MTU minus 20 bytes of IP header and 8 bytes of UDP header).
A character takes 1 to 3 bytes in UTF-8, so 500 characters make 500 to 1\,500 bytes of text.
Measured with the sender name ``Ana'' (Figure~\ref{fig:sizes}):

\begin{center}
\small
\begin{tabular}{lcc}
\toprule
500 characters of & general message & private message, encrypted \\
\midrule
ASCII (1 byte per letter)                       & 590 bytes   & 819 bytes \\
Romanian or Cyrillic (2 bytes per letter, as \u{a}) & 1\,090 bytes & 1\,487 bytes \\
Chinese (3 bytes per character)                 & 1\,590 bytes & 2\,155 bytes \\
\bottomrule
\end{tabular}
\end{center}

Ordinary messages are far below one frame; only the longest messages in non-Latin alphabets need
two IP fragments, which a local network handles well. 500 characters is a limit at which a message
is always one datagram of reasonable size.

A message of 100\,000 characters would be at least 100\,000 bytes of text, more than the largest UDP
datagram of IPv4 (65\,507 bytes of data, Section~\ref{sec:control}, question~3): \code{sendto}
would fail with ``message too long''. Even a message just below that limit would be cut by IP into
about 45 fragments; one lost fragment loses the whole datagram, and every fragment is a chance to
lose it. To send it anyway, the program would have to cut it into numbered parts and put them
together again at the receiver, which is a part of TCP built again. Our chat refuses a line of more
than 500 characters before sending it (\code{* Too long: 612 characters, at most 500. Not sent.}),
and drops any received datagram longer than 4\,096 bytes or with a longer text.

\begin{figure}[H]
  \centering
  \fig{datagram_sizes.png}{\linewidth}
  \caption{The size of every kind of datagram, and of the longest messages, against the payload of
           one Ethernet frame (1\,472 bytes) and the limit above which the chat drops a datagram.}
  \label{fig:sizes}
\end{figure}

% Section 4
\section{Task 2: Who is online}

Every peer keeps a table of the peers online: name, id, endpoint of the own socket, and the time it
was last heard from. The protocol that fills it:

\begin{lstlisting}
new peer N                      the peers online (A, B)            the peer that has N's name
 |-- hello  (to the group) -------->|                                  |
 |<-- here  (A, to N's own port) ---|  at once                         |
 |<-- here  (B, to N's own port) ---|                                  |
 |<-- taken (to N's own port) ---------------------------------------- | if the name is used
 |   ... waits 1 second; with taken: prints that the name is used, ends
 |-- join   (to the group) -------->|  A and B add N: "* N joined the chat"
 |-- alive  (to the group, every 5 s) ...
 |-- bye    (to the group, on /quit) -->|  "* N left the chat"
\end{lstlisting}

The others add the new peer only when it says \code{join}, after the second in which the owner of
its name could object. A peer that is refused therefore never appears in anybody's list. A
\code{join} replaces an older peer with the same name, because that one is gone (killed and
started again at once), or it would have answered \code{taken}; the test \code{restarted peer}
checks that a private message then reaches the new process. Any datagram of a peer counts as a sign
of life. An \code{alive} or a message from a peer that is not in the table adds it, if its name is
free: its \code{join} was lost, or it joined while we were joining ourselves.

\subsection{Answers to the questions of Task 2}

\question{1. How does a new peer learn who is already online, and how long does it take?}

It sends \code{hello} to the group. Every peer that is online answers at once with \code{here},
sent to the source of the \code{hello}, which is the own socket of the new peer; the answer carries
the name and id of the peer, and its source is the endpoint of that peer. The new peer waits one
second, collecting the answers, and then knows who is online with names and endpoints
(\code{* You joined the chat as Maria. Online: Ana, Ion.}); \code{/users} lists them:

\begin{lstlisting}
* 3 online:
    Maria (you)            192.168.100.15:63791
    Ana                    192.168.100.15:60685   last heard 0 s ago
    Ion                    192.168.100.15:60686   last heard 0 s ago
\end{lstlisting}

Measured over ten joins with three peers online, the last answer arrived \textbf{21\,ms} after the
\code{hello} (median, from the log of the new peer). Nearly all of that is the new process itself:
it compiles its receiving code at the first datagram (the .NET just-in-time compiler), while a
private message between two running peers is acknowledged in 0.6\,ms (Task~3). On a real Wi-Fi
network, a few milliseconds more. The one-second wait is therefore generous; it is there for a
slow peer or a busy network, and it is the time the owner of a name has to object.

\question{2. The process of a peer is killed. How and when do the others notice? Why can they not
notice it at once, as a TCP server notices a closed connection?}

A killed process sends nothing more: no \code{bye}, no \code{alive}. The others notice the silence:
the network thread removes a peer that has not been heard from for 15 seconds and prints
\code{* Ion left the chat (not heard from for 15 s)}. Since the last \code{alive} was sent 0 to 5
seconds before the kill, the others notice it \textbf{10 to 15 seconds} after the kill (plus up to
0.2\,s until the next check). Eight kills at random moments were noticed after 10.2 to 14.3\,s:

\begin{figure}[H]
  \centering
  \fig{killed_peer.png}{\linewidth}
  \caption{Left: the 15 seconds start at the last datagram, not at the kill. Right: the measured
           delay between the kill and the moment the other peer showed it, for eight kills.}
\end{figure}

A TCP server notices at once because there is a connection: when a process ends, its operating
system closes its TCP sockets and sends a FIN (or a RST) to the other side, whose \code{recv} then
returns 0 (or fails). UDP has no connection and no state shared by the two sides: closing a UDP
socket sends nothing, and the peers never agreed that they are talking. Only the absence of
datagrams can show that a peer is gone, and silence takes time to be sure about: 15 seconds are
three missed \code{alive} messages, so one or two lost datagrams do not remove a peer that is still
there. (A TCP server would not notice at once either if the network of the client disappeared;
laboratory work~2 needed keepalive and took 16\,s for it.)

\question{3. Two peers start with the same name at the same moment. What happens?}

Both send \code{hello}, and both are still in their second of joining when they receive the
other's \code{hello} with their own name. The rule is the same in both: \textbf{the smaller id keeps
the name}. The peer with the greater id stops joining and prints
\code{* Another peer is joining with the name Bob at this moment. Start again with another -{}-name.};
the one with the smaller id answers \code{taken} to the other and joins normally. The peers that
are already online answer both with \code{here}, but they add a peer only on its \code{join}, which
only the winner sends, so they see one \code{* Bob joined the chat}. The test
\code{same name same moment} starts two peers with the same name right after each other, three
times: each time exactly one joined, and the observer saw it join once.

If the \code{hello} of both were lost, both would join. Then, as soon as one of them hears the other
(\code{join}, \code{alive} or a message with its own name and another id), the same rule applies:
the one with the greater id says \code{bye} and ends, within 5 seconds at most.

% Section 5
\section{Task 3: Private messages}

\code{/msg Ion see you at five} looks Ion up in the table and sends one \code{private} datagram from
the own socket of the sender to the own socket of Ion:

\begin{lstlisting}
{"proto":"lanchat/1","type":"private","from":"Ana","id":"69478b4234df7e63","seq":2,"to":"Ion",
 "text":"see you at five"}
\end{lstlisting}

The sender shows \code{[private to Ion] Ana: see you at five}, the receiver
\code{[private] Ana: see you at five}. A name that is not online gets
\code{* Nobody is not online. /users shows who is.}, and nothing is sent. The receiver drops a
private message that arrives on its shared socket, or whose \code{to} is not its name.

The receiver answers every private message with an \code{ack} carrying its sequence number. If the
\code{ack} has not arrived after one second, the sender sends \emph{the same datagram} again, at
most three times in all; then it tells the user
\code{* Ion did not confirm your private message "see you at five": it was probably not delivered.}
The receiver acknowledges a copy too, because a copy means that its first \code{ack} was lost, but
shows it only once (it remembers the pairs id and sequence number).

The test \code{private message} runs three peers and, on the same computer, a stranger: a plain
socket bound to the port of the chat and joined to the group. After Ana's private message to Ion,
Ion shows it, Maria does not, the stranger received nothing private, Ana's log has exactly one
\code{private} datagram sent, to Ion's own endpoint, and one \code{ack} back.

\subsection{Answers to the questions of Task 3}

\question{1. Several peers on one computer share the port of the chat. How does a private message
reach the right one?}

It is not sent to the shared port. Every peer has a second socket, bound to port 0, for which the
system chooses a free port that only this peer uses (in the session of Section~2.4, \code{60686}
for Ion and \code{63791} for Maria, both on \code{192.168.100.15}). Every datagram of a peer leaves
from that socket, so its source endpoint is the address of the peer, and the others keep it in their
table. Ana's private message went to \code{192.168.100.15:60686}, where only Ion's socket is bound: the
system delivers it to that socket and to no other. A unicast datagram to the shared port would have
reached only the first of the sockets bound to it (Section~\ref{sec:task1}, question~2), which is
Ion only by chance.

\question{2. Why is a private message not sent to everybody, with the name of its receiver inside?}

Because then everybody would receive it, and the privacy would depend on every other program
politely ignoring a message that is not addressed to it. Any program bound to the port, a modified
chat, or Wireshark on the same network would read it. The cost would also be everybody's: every
peer would receive, parse and drop every private message, and on Wi-Fi a multicast datagram is sent
at the lowest data rate and without the acknowledgements and retransmissions of the Wi-Fi layer that
a unicast frame has, so it is slower and is lost more often. Sent to one endpoint, the message
reaches only the receiver's computer and socket (on a switched network not even the network cards of
the others see it), and the Wi-Fi layer retransmits it if a frame is lost.

\question{3. What happens to a private message sent to a peer that left a second ago?}

It depends on how the peer left. If it typed \code{/quit}, it sent \code{bye}, it is no longer in
the table, and the message is not sent at all: \code{* Ion is not online}.

If its process was killed, the others still have it in their table for up to 15 seconds, and the
message is sent to its old endpoint. The datagram arrives at its computer, but no socket is bound to
that port any more: the computer answers with an ICMP ``port unreachable'' message. \code{sendto}
has already succeeded and knows nothing about it. On Windows, the next \code{recvfrom} on the
sending socket fails with \code{WSAECONNRESET} (error 10054); our chat catches it, looks at the last
datagram it sent to one peer, and tells the user at once (test \code{private to gone peer}):

\begin{lstlisting}
* Ion did not get your private message: nobody listens at 192.168.100.15:60003 any more.
00:15:08.896  ----  ICMP port unreachable: nobody listens at 192.168.100.15:60003 any more, where the last private went
\end{lstlisting}

The error does not say which datagram it is about, so the chat can only guess that it is the last
one. On Linux an unconnected UDP socket does not report it at all, and when the receiving computer is
gone, or its firewall drops the datagram without an answer, nothing comes back either. Then the
\code{ack} does not come, the message is sent three times, and the sender is told that it was
probably not delivered. Either way the message is lost: UDP does not keep it for a peer that comes
back.

\question{4. UDP can lose a private message. How could the sender know that it arrived, and what
would it do when it did not?}

Only the receiver can tell, so it must answer: our chat implements this. The receiver sends an
\code{ack} with the sequence number of every private message it accepts. The sender keeps the
message until the \code{ack} arrives; after 1 second without it, it sends the same datagram again,
at most three times in all, and then tells the user that the message was probably not delivered.
The receiver acknowledges every copy but shows a message only once, recognising copies by the pair
(id, sequence number). The time to wait must be longer than a round trip: measured on one computer,
an \code{ack} came back after 0.6\,ms (median of 100, timed by a test program), so one second
leaves a wide margin for a real network.

Two limits remain. If the message arrived but every \code{ack} was lost, the sender wrongly reports
it as not delivered: the sender can never be sure, only very likely right. And the messages are not
ordered: a message sent again may be shown after a later one. TCP adds sequence numbers for the
order and a window to send many messages before the first acknowledgement; one message at a time
does not need them.

% Section 6
\section{Task 4: Secure private messages}

All the peers of a chat are started with the same \code{-{}-passphrase}. Each derives the same key
from it, without exchanging anything, and every private message is encrypted with it; the general
channel stays in clear (\code{PrivateCipher.cs}):

\begin{center}
\small
\begin{tabularx}{\linewidth}{@{}lX@{}}
\toprule
Step & What is used \\
\midrule
Key derivation & PBKDF2 with HMAC-SHA256, 600\,000 iterations, salt \code{"lanchat/1 port 5000"},
  32 bytes of key \\
Encryption & AES-256-GCM, with a 16-byte authentication tag \\
Nonce & 12 random bytes from the system's cryptographic generator, new for every message \\
Also authenticated & \code{lanchat/1|from|id|to|seq}, the header of the message, as associated data \\
\bottomrule
\end{tabularx}
\end{center}

The encrypted private message has no \code{text}; \code{data} holds the ciphertext followed by the
tag (seen by a test program that joined the chat without the passphrase):

\begin{lstlisting}
{"proto":"lanchat/1","type":"private","from":"Ana","id":"3e98cd0b4749bd55","seq":2,"to":"Eve",
 "nonce":"DvQri0iAX/4B8nWI","data":"VTGb6epuIfsfnQAIcxidl..."}
\end{lstlisting}

A peer with another passphrase computes another key: the tag does not match, the message is dropped
(\code{cannot be decrypted: another passphrase, or the datagram was changed on the way}), no
\code{ack} is sent, and the sender is told that it was not confirmed (test
\code{wrong passphrase}). A peer without a passphrase drops an encrypted message, and a peer with
one drops a message that is not encrypted.

\subsection{Answers to the questions of Task 4}

\question{1. Why is the key derived with a key-derivation function, and not taken from the bytes of
the passphrase?}

First, AES-256 needs exactly 32 bytes, and a passphrase has any length: its bytes would have to be
cut or padded in some arbitrary way. Second, and more important, a key should be 32 bytes that are
completely unpredictable, while a passphrase chosen by people is a few words or a name with digits,
which an attacker can guess. Whoever captures one encrypted datagram can test a guess offline: derive
a key and check whether the tag of the datagram matches. With the bytes of the passphrase as the key,
one guess costs one AES operation, and billions of guesses per second are possible.

A key-derivation function makes every guess expensive on purpose. PBKDF2 runs HMAC-SHA256 600\,000
times (the number OWASP recommends) to produce the key. Deriving our key took 71 to 75\,ms in the C\#
peer and 0.3 to 0.5\,s in the Java one; a peer pays that once at start, an attacker for every single
guess. The salt makes the result depend on more than the passphrase, so that a table computed in
advance for common passphrases does not apply. Usually the salt is random and stored next to the
key; here every peer must derive the same key without exchanging anything, so the salt is fixed by
the protocol and the port of the chat. Two chats on different ports with the same passphrase still
have different keys; a precomputed table for our protocol and port 5000 remains possible, which a
strong passphrase makes useless.

\question{2. Which algorithm encrypts your messages? Why must its nonce (or IV) never be used twice
with the same key, and where do yours come from?}

AES-256 in Galois/Counter Mode (GCM), an authenticated encryption: it encrypts and, in the same
step, computes a 16-byte tag over the ciphertext and the associated data. GCM encrypts by counter
mode: AES encrypts the nonce followed by a counter, and the result, the keystream, is XORed with the
text. The keystream depends only on the key and the nonce. If two messages used the same nonce, the
XOR of their ciphertexts would be the XOR of their texts, from which the texts can often be
recovered, and the authentication key of GHASH could be computed, which allows forging tags for any
message. One repeated nonce breaks both the secrecy and the integrity.

Our nonces are 12 bytes, new for every message, from the cryptographic random generator of the
system: \code{RandomNumberGenerator} in C\# and \code{SecureRandom} in Java. There are
$2^{96}$ possible nonces; the chance of a repetition among $n$ messages is about $n^2/2^{97}$, which
is negligible far beyond any chat (NIST allows $2^{32}$ messages per key with random nonces). A
counter would not be safe here without more care, because all the peers share the key: two peers
counting from 1 would use the same nonces. A message sent again because its \code{ack} was missing
is the same datagram, with the same nonce, ciphertext and tag: the same text with the same nonce
reveals nothing new.

\question{3. Every member of the chat knows the passphrase. Who can read a private message between
Ana and Ion? What would be needed so that only the two of them can?}

Every member. The key comes only from the passphrase, so Maria has the same key as Ana and Ion. The
datagram is sent to Ion's socket only, but on a Wi-Fi network Maria can capture it (a network
card in monitor mode, Wireshark, or an access point she controls), and her copy of the key decrypts
it. The encryption protects the private messages from those who do not know the passphrase, not from
the other members.

For only Ana and Ion to read it, each pair of peers needs a key of its own, which the others cannot
compute: a key exchange. Each peer generates a key pair at start (for example X25519) and announces
its public key; Ana and Ion each combine their own private key with the other's public key
(Diffie-Hellman) and obtain the same secret, which nobody who only saw the public keys can compute;
from it a key for the pair is derived (with HKDF) and used for their messages. The public keys must
also be authenticated, otherwise Maria could announce her own key as Ion's and read everything in
the middle: by comparing a fingerprint of the key in person, or with keys signed by an authority.
This is how Signal or TLS work. Changing the key from time to time would also protect old messages
if a key leaks later (forward secrecy).

\question{4. Can somebody who does not know the passphrase change an encrypted message without the
receiver noticing? What protects against it?}

No. The GCM tag is computed with the key over the ciphertext and over the associated data, which is
the header of the message: protocol, sender, id, receiver and sequence number. Changing any bit of
any of them makes the receiver compute another tag, and without the key nobody can compute the right
one (the chance of guessing a 16-byte tag is $2^{-128}$). The receiver then drops the message,
showing nothing of it. The test \code{tampered message} captured a real encrypted datagram and sent
changed versions of it to a peer with the passphrase:

\begin{center}
\small
\begin{tabularx}{\linewidth}{@{}lX@{}}
\toprule
Change & What the receiver did \\
\midrule
one bit of the ciphertext flipped & dropped: cannot be decrypted \\
another sequence number (header) & dropped: cannot be decrypted \\
another nonce & dropped: cannot be decrypted \\
\code{text} in clear instead of \code{data} & dropped: not encrypted, but this chat uses a passphrase \\
the same datagram, unchanged, to a peer that never saw it & \textbf{accepted}, then dropped as a copy \\
\bottomrule
\end{tabularx}
\end{center}

The last row shows what the tag does not protect against: a \emph{replay}. An unchanged datagram is
authentic, so it is accepted. Our receiver remembers the (id, sequence number) of the messages it has
shown, so a replay to the same running peer is dropped as a copy; but a peer started again with the
same name has an empty memory and accepted it once. Protecting against that needs the receiver to
be part of the key or of the authenticated data in a way that changes with every session (a key per
pair and per session, as in question~3, or timestamps). The general channel and the control
messages (\code{join}, \code{bye}, \code{taken}) are not authenticated at all: anybody on the
network can send a general message in Ana's name, or tell a peer that its name is taken. The task
asked for private messages to be secret; authenticating the rest would need a tag on every datagram.

% Section 7
\section{Control questions}\label{sec:control}

\question{1. Which promises of TCP does UDP not make? Which of them did your chat rebuild, and which
did it do without?}

TCP promises a connection (both sides agree to talk, and each learns when the other closes), reliable
delivery (acknowledgements and retransmission), the original order, no duplicates, flow control (the
sender never overruns the receiver's buffer) and congestion control (it slows down when the network
is loaded). UDP promises none of them: a datagram may be lost, duplicated or reordered, and nobody is
told. It only keeps each datagram whole, which TCP does not.

The chat rebuilt three of them, in a small form. Presence instead of a connection: \code{hello},
\code{join}, \code{alive} every 5\,s, \code{bye}, and 15\,s of silence for a peer that is gone.
Reliable delivery for private messages: \code{ack}, three attempts, one second apart. No duplicates:
the (id, sequence number) of the last 1\,000 messages. It did without the order (two quick messages
may be shown in another order, especially when one is sent again), without reliable delivery for
the general channel (a lost multicast datagram is simply missing for that peer; acknowledgements
from every peer would multiply the traffic), and without flow and congestion control (a chat sends
a few small datagrams per minute).

\question{2. Why is one sendto always one recvfrom with UDP, while one send is not one receive with
TCP? Why did your chat need no delimiter and no length prefix?}

UDP is a protocol of messages: every \code{sendto} becomes one datagram with its own UDP header,
which contains its length, and the receiving system delivers it whole to one \code{recvfrom}, or not
at all; it never joins two datagrams or delivers half of one. TCP is a stream of bytes: \code{send}
only copies bytes into the sending buffer of the system, which cuts them into segments as the network
allows (maximum segment size, window, Nagle's algorithm), and \code{recv} returns whatever bytes have
arrived, part of a message or several messages. That is why laboratory work~2 needed the \code{\textbackslash n}
at the end of every line and a class that collects bytes until a line is complete. In the chat, one
datagram is one JSON object: its end is the end of the datagram, and its length is in the UDP header.

\question{3. How big can a UDP datagram be? Why should it stay below about 1\,400 bytes on a real
network?}

The length field of the UDP header has 16 bits, so a UDP datagram, header included, has at most
65\,535 bytes; inside an IPv4 packet, whose total length is also at most 65\,535 bytes, that leaves
$65\,535 - 20 - 8 = 65\,507$ bytes of data. But one Ethernet or Wi-Fi frame carries at most 1\,500
bytes of IP packet (the MTU), so 1\,472 bytes of UDP data. A larger datagram is cut by IP into
fragments, which are reassembled only at the receiver. The whole datagram is lost if any fragment is
lost, so the risk grows with every fragment; some firewalls and NAT devices drop fragments
altogether; and on many paths the MTU is smaller than 1\,500 (1\,492 with PPPoE, less in a VPN
tunnel), so a datagram of 1\,472 bytes may already be fragmented there. About 1\,400 bytes leaves a
margin for such extra headers. The datagrams of our chat are 70 to 170 bytes for everything except
long messages; 500 Latin characters make 590 bytes, and only the longest messages in non-Latin
alphabets are over one frame (Figure~\ref{fig:sizes}).

\question{4. What is the difference between broadcast and multicast? Which one does your chat use
for the general channel, and what would change with the other one?}

A broadcast goes to every host of the network segment: to \code{255.255.255.255} (limited broadcast)
or to the broadcast address of the network, here \code{192.168.100.255}. Every network card accepts
it, and the operating system of every computer must look at it, whether or not a program wants it;
routers do not pass it; IPv6 has no broadcast. A multicast datagram goes to a group address
(\code{224.0.0.0/4}). Only the hosts that have joined the group receive it: joining tells the network
card to accept the Ethernet address of the group and, through IGMP, tells switches with IGMP snooping
to send the group only to the ports that want it. With a larger time to live and multicast routing,
it can even cross routers.

Our chat uses \textbf{multicast} to \code{239.255.42.42} by default, and broadcast with
\code{-{}-broadcast}. With broadcast, the sending socket needs \code{SO\_BROADCAST} (without it
\code{sendto} fails), the destination is the broadcast address of the network, computed from the
address and the mask of the adapter, and the receivers need not join anything. Every computer on the
Wi-Fi network then receives every message of the chat, even those that do not run it. In exchange,
broadcast works on networks whose switches handle IGMP badly. Because every shared socket is bound to
\code{0.0.0.0} and has joined the group, it receives both, so broadcast and multicast peers can be in
the same chat (test \code{broadcast mode}).

\question{5. How many sockets does your chat use? To which address and port is each of them bound,
and why?}

Two per peer (Section~2.1). The shared socket is bound to \code{0.0.0.0} and the port of the chat,
with \code{SO\_REUSEADDR}, and joins the group: the general channel goes to that port, which all the
peers know in advance and share, also on one computer, and \code{0.0.0.0} accepts datagrams addressed
to the group or to the broadcast address. The own socket is bound to \code{0.0.0.0} and port 0, a free
port chosen by the system and used by this peer only: a unicast datagram to the shared port would
reach only one of the peers of a computer, so private messages, answers and acknowledgements need a
port of their own. The own socket sends every datagram, so its endpoint is the address of the peer,
which the others learn from the source of what it sends. The Java peer has the same two sockets.

\question{6. What does the checksum in the UDP header protect against? Why is it not enough against
somebody who changes your datagrams on purpose?}

The checksum is a 16-bit ones' complement sum over the UDP header, the data, and a pseudo-header with
the source and destination addresses, the protocol and the length. The receiving system recomputes it
and drops a datagram whose sum does not match. It protects against accidental changes: bits damaged
in transit or in the memory of a device on the way, and a datagram delivered to the wrong address. It
is weak even for that (some combinations of errors cancel out), and in IPv4 it is optional.

Against an attacker it protects nothing, because it uses no secret: anybody who changes a datagram
simply computes the checksum of the new content, with the same public algorithm, and the receiver
accepts it. Protection against deliberate changes needs a value that only those who know a key can
compute: a message authentication code, like the GCM tag of our private messages (Task~4,
question~4).

% Section 8
\section{Bonus: a chat in two languages}

The Java peer (\code{lab3/chat-java}, 1\,200 lines) implements the same protocol, from
\code{PROTOCOL.md}, independently of the C\# code. It joins the same chat, finds the C\# peers and is
found by them, shows their general messages, and exchanges encrypted private messages with them in
both directions. The differences are those of the language:

\begin{itemize}
  \item \code{DatagramSocket} has no \code{select}, so the Java peer has a thread for each socket and
        one for the timers; everything they share is protected by one lock.
  \item Java has no JSON in its standard library, so \code{Json.java} writes and reads it.
  \item Java cannot list the gateways of the adapters. To find the network, the Java peer
        ``connects'' a UDP socket to an address outside the local network: that sends nothing, but
        makes the system choose the route, and with it the local address of the default adapter.
  \item \code{PBKDF2WithHmacSHA256} and \code{AES/GCM/NoPadding} give exactly the same key and
        ciphertext format as \code{Rfc2898DeriveBytes.Pbkdf2} and \code{AesGcm}: the Java cipher puts
        the tag after the ciphertext, and the C\# peer does the same.
\end{itemize}

The test \code{java peer} starts a C\# peer, a Java peer and another C\# peer with one passphrase
containing Romanian letters, and checks general messages in both directions, private messages from
C\# to Java and from Java to C\#, the acknowledgements, \code{/users} in Java, that a Java peer cannot
take the name of a C\# one, and the \code{bye} of the Java peer. The test \code{java invalid
datagrams} sends the sixteen invalid datagrams of Section~\ref{sec:tests} to the Java peer: it drops
each one, with the same reason in its log as the C\# peer.

\notebox{\textbf{A difference found by the tests.} With the passphrase given on the command line, the
Java peer could not read the C\# peer's private messages. On Windows the Java launcher decodes the
command line with the ANSI code page of the system (here Cp1251), which has no letter \u{a}, so it
arrived as a question mark and gave another key; the C\# runtime reads the command line as Unicode.
The environment is read as Unicode by both, so the Java peer also accepts the passphrase from the
environment variable \code{LANCHAT\_PASSPHRASE}, and warns when a passphrase on the command line
contains a question mark.}

% Section 9
\section{Tests}\label{sec:tests}

The tests (\code{lab3/tests/test\_chat.py}) start real peers on this computer, each on a port of its
own, type into them through their standard input, and check what they print and what their logs say.
Where a test needs to send something the chat never sends, it uses plain Python sockets. All 19 tests
passed:

\begin{center}
\small
\begin{tabularx}{\linewidth}{@{}lXc@{}}
\toprule
Test & What it checks & Result \\
\midrule
general channel & three peers; each message shown by all, the author too, once; one datagram to the group & passed \\
join and users & the newcomer knows the two peers online at once; \code{/users} shows their endpoints & passed \\
name taken & a second ``ana'' is told that the name is used and ends (exit code 3) & passed \\
same name same moment & three times two peers with one name at once: exactly one joins & passed \\
bye & \code{/quit}: the others show it 6\,ms later & passed \\
killed peer & a killed peer is shown as gone 10 to 15\,s later & passed \\
restarted peer & a peer killed and started again replaces its old self & passed \\
private message & only the receiver gets it, not a peer or a stranger on the same computer; one \code{ack} & passed \\
private errors & name not online, missing text, unknown command, 501 characters refused, 500 sent & passed \\
private to gone peer & the ICMP ``port unreachable'' is reported to the user & passed \\
invalid datagrams & 16 kinds of invalid datagrams dropped, each with its reason in the log; the chat goes on & passed \\
unicode & names and messages in Romanian and Russian; names not case sensitive & passed \\
encrypted on the wire & a stranger in the chat receives \code{nonce} and \code{data}, no text & passed \\
wrong passphrase & a peer with another passphrase, or none, shows nothing; three attempts reach it & passed \\
tampered message & changed datagrams dropped; an unchanged replay accepted once & passed \\
broadcast mode & broadcast and multicast peers in one chat & passed \\
loopback only & a chat on \code{127.0.0.1}, between the peers of one computer & passed \\
java peer & C\# and Java peers in one encrypted chat (bonus) & passed \\
java invalid datagrams & the Java peer drops the same invalid datagrams & passed \\
\bottomrule
\end{tabularx}
\end{center}

The invalid datagrams were: 300 random bytes, a sentence, a JSON array, JSON without the protocol, an
unknown kind, a message without text, a text of 501 characters, a name with a space, a message without
a sequence number, datagrams of 5\,000 and 60\,000 bytes, a private message without text, a private
message sent to the shared port, a private message for somebody else, an \code{ack} nobody waits for,
and a \code{bye} from a peer that is not online.

The interactive console was checked by hand: a peer ran in a terminal while another one sent it
messages; the messages appeared above the prompt and the prompt stayed below them. Ctrl+C in that
terminal sent \code{bye} before the program ended.

\subsection{On two computers}\label{sec:two}

\notebox{The assignment requires at least three peers on at least two computers. Everything in this
report ran on one computer: the datagrams of the general channel went out through the Wi-Fi adapter,
but no datagram came from another computer. Run Task~1 to 4 on two computers (the commands are in
Appendix~\ref{app:run} and in \code{lab3/README.md}) and fill in the result here before the defense.}

\begin{center}
\small
\begin{tabularx}{\linewidth}{@{}p{5.6cm}X@{}}
\toprule
Network & \TODO{e.g. hotspot of a phone, 192.168.43.0/24} \\
Computer 1 (C\# peers) & \TODO{address, peers} \\
Computer 2 (C\# and Java peers) & \TODO{address, peers} \\
General channel, private messages, encryption & \TODO{works / what did not} \\
\bottomrule
\end{tabularx}
\end{center}

% Section 10
\section{Conclusion}

UDP gives a program almost nothing but the delivery of separate messages, most of the time. That is
enough for a chat without a server, because what the chat needs from the network is exactly what UDP
does well: one datagram to the multicast group reaches every peer of the network at once, and one
datagram to an endpoint reaches one peer. Everything else had to be built on top, and each part
answers one of the questions of the assignment:

\begin{center}
\small
\begin{tabularx}{\linewidth}{@{}p{4.4cm}X@{}}
\toprule
Need & How the chat meets it \\
\midrule
everybody on one port & \code{SO\_REUSEADDR}, \code{0.0.0.0}, a multicast group with TTL 1 \\
one peer only & a second socket on a port of its own, which sends everything \\
recognise its own message & a random id in every datagram \\
who is online & \code{hello} answered by \code{here}, \code{join}; measured: all answers within 21\,ms \\
who has left & \code{bye}; or 15\,s of silence, noticed 10 to 15\,s after a kill \\
unique names & \code{taken}; the smaller id wins a tie \\
lost private messages & \code{ack}, three attempts, one second apart; copies shown once \\
secret private messages & PBKDF2 (600\,000 iterations) and AES-256-GCM with random nonces \\
changed messages & the GCM tag over the text and the header \\
invalid datagrams & checked field by field, dropped and logged \\
\bottomrule
\end{tabularx}
\end{center}

The tests also showed the limits honestly. A peer that is killed cannot be noticed at once, only
after a silence that is long enough to be sure. A sender can never be completely sure that a message
arrived. A passphrase shared by all the members protects the private messages from strangers, not
from the other members, and authentication alone does not stop a replayed message. The Java peer,
written from the protocol description, talked to the C\# one without changes to either, apart from
the code page of the Java command line on Windows.

% Appendices
\appendix

\section{Building and running}\label{app:run}

The C\# chat needs the .NET 9 SDK, the Java peer a JDK from version 17, the tests Python 3. The
commands run from the repository root, every peer in its own terminal.

\begin{lstlisting}
# build
dotnet build Labs.csproj -c Release
javac -d lab3/chat-java/out lab3/chat-java/src/lanchat/*.java

# peers (same port, and same passphrase if any, for every peer of one chat)
bin\Release\net9.0\Labs.exe lab3 chat --name Ana --port 5000
bin\Release\net9.0\Labs.exe lab3 chat --name Ion --port 5000 --passphrase "our secret"
bin\Release\net9.0\Labs.exe lab3 chat --name Maria --port 5000 --broadcast
bin\Release\net9.0\Labs.exe lab3 chat --name Ion --port 5000 --interface 192.168.43.12
bin\Release\net9.0\Labs.exe lab3 networks           # the addresses for --interface
java -cp lab3/chat-java/out lanchat.LanChat --name Dan --port 5000 --passphrase "our secret"

# in the chat
hello everyone                 # to everybody
/msg Ion see you at five       # to Ion only
/users                         # who is online
/quit                          # leave (also Ctrl+C)

# tests and measurements
python lab3/tests/test_chat.py
python lab3/tests/test_chat.py measure      # writes lab3/results/measurements.json
python lab3/report/make_charts.py           # the figures of this report
\end{lstlisting}

On two computers, the network must let multicast (or, with \code{-{}-broadcast}, broadcast) through;
many Wi-Fi networks do not, and the hotspot of a phone does. The first time, Windows asks whether
\code{Labs.exe} and \code{java.exe} may use the network: they must be allowed on private networks.

\section{Measurements}\label{app:data}

The measurements were made with \code{python lab3/tests/test\_chat.py measure} and saved in
\code{lab3/results/measurements.json}, from which \code{make\_charts.py} draws the figures. Timings
from the log files have steps of 1\,ms.

\begin{center}
\small
\begin{tabularx}{\linewidth}{@{}Xl@{}}
\toprule
Measurement & Result \\
\midrule
last \code{here} after \code{hello}, three peers online (10 joins) & median 21\,ms, at most 24\,ms \\
\code{private} to \code{ack}, timed by a test program (100 messages), C\# peer & median 0.62\,ms \\
the same, Java peer & median 1.11\,ms \\
\code{private} to \code{ack} in the logs, encrypted or not, C\# and Java (20 each) & median 0 or 1\,ms, at most 12\,ms \\
key derivation, 600\,000 iterations, C\# / Java & 71 to 75\,ms / 0.3 to 0.5\,s \\
killed peer noticed after (8 kills) & 10.2 to 14.3\,s \\
\code{hello} / \code{join} / \code{here} / \code{ack} (name ``Ana'') & 73 / 72 / 72 / 79 bytes \\
private message of 9 or 10 characters, not encrypted / encrypted & 113--115 / 167--168 bytes \\
\bottomrule
\end{tabularx}
\end{center}

\end{document}
